% Wersja źródłowa: audyt II, jednostka 22

\section{Dystrybucje temperowane i transformacja Fouriera}
\label{sec:temperowane22}

\providecommand{\atlasref}[2]{#2}


\noindent\textbf{Fakty z wcześniejszych rozdziałów.}
Przyjmujemy, że przestrzeń Schwartza jest zupełną przestrzenią Frécheta,
$\mathcal D(\R^d)$ jest w niej gęsta, a transformacja Fouriera
jest jej ciągłym automorfizmem z odwrotnością i wzorem
$\mathcal F^2\varphi=\varphi(-\cdot)$.
Dystrybucja o zwartym nośniku $K$ działa również na funkcjach
$C^\infty$ poprzez $\langle T,f\rangle=\langle T,\chi f\rangle$,
gdzie $\chi=1$ w otoczeniu $K$, i spełnia oszacowanie
skończonego rzędu na ustalonym zwartym otoczeniu $K$.
Korzystamy z Twierdzeń~\ref{thm:schwartz-frechet17},
\ref{thm:density-schwartz17}, \ref{thm:fourier-schwartz17}
i Lematu~\ref{lem:compact-action21}. Wzór na $\mathcal F^2$
udowodniono we Wniosku~\ref{cor:dt-fourier-iteration}.

W całym rozdziale $d\ge1$ jest liczbą całkowitą.
Rozważamy funkcje oraz dystrybucje zespolone na $\R^d$,
chyba że jawnie dopuszczamy także ciało $\R$.
Przestrzeń Schwartza $\mathcal S=\mathcal S(\R^d;\C)$
wyposażamy w topologię Frécheta z rozdziału~\ref{sec:gladkie-schwartz17}.
Dla $N,m\in\N_0$ przyjmujemy seminormy (również oznaczane
w dalszej części przez $p_{N,m}=Q_{N,m}$)
$Q_{N,m}(\varphi)=\max_{|\alpha|\le m}\sup_{x\in\R^d}
(1+|x|)^N|\partial^\alpha\varphi(x)|$.
Przez $\mathcal S'$ rozumiemy przestrzeń \emph{ciągłych,
$\C$-liniowych} funkcjonałów na $\mathcal S$.
W parowaniu $\langle T,\varphi\rangle$ nie sprzęgamy zespolenie
funkcji testowej. Przypominamy, że
$\mathcal D(\R^d)=C_c^\infty(\R^d)$ jest gęsta w $\mathcal S$.

\subsection{Dystrybucje temperowane jako przestrzeń sprzężona}

\noindent\emph{Powiązania tematyczne w Bibliotece matematycznej:} \atlasref{PDE-08::def:pde8-tempered}{PDE-08: definicja dystrybucji temperowanej}.


\begin{definicja}[Dystrybucja temperowana]
\label{def:tempered22}
Dystrybucją temperowaną na $\R^d$ nazywamy element
$T\in\mathcal S'$, czyli ciągły liniowy funkcjonał
$T:\mathcal S\to\C$. Rozpatrujemy na $\mathcal S'$ topologię
słabą-* $\sigma(\mathcal S',\mathcal S)$, wyznaczoną przez
seminormy $T\mapsto|\langle T,\varphi\rangle|$
($\varphi\in\mathcal S$), oraz topologię silną
$\beta(\mathcal S',\mathcal S)$, wyznaczoną przez
$P_B(T)=\sup_{\varphi\in B}|\langle T,\varphi\rangle|$,
gdzie $B$ przebiega wszystkie niepuste zbiory ograniczone w $\mathcal S$.
Dla każdego ciągłego funkcjonału $T$ liczba $P_B(T)$
jest skończona, ponieważ ciągłe odwzorowania liniowe
przenoszą zbiory ograniczone na zbiory ograniczone.
\end{definicja}

\begin{twierdzenie}[Kryterium wzrostu seminormowego]
\label{thm:tempered-estimate22}
Niech $T:\mathcal S(\R^d;\C)\to\C$ będzie liniowy.
Następujące warunki są równoważne:
\begin{enumerate}
\item $T\in\mathcal S'(\R^d)$;
\item istnieją $C>0$ oraz $N,m\in\N_0$ takie, że
$|\langle T,\varphi\rangle|\le C Q_{N,m}(\varphi)$
dla wszystkich $\varphi\in\mathcal S$.
\end{enumerate}
W szczególności ograniczenie $T|_{\mathcal D}$ jest dystrybucją,
a odwzorowanie $\mathcal S'\to\mathcal D'$ dane przez
ograniczenie jest liniowe i różnowartościowe.
\end{twierdzenie}
\begin{proof}
Rodzina $Q_{j,j}$, $j\in\N_0$, jest rosnąca
i wyznacza topologię $\mathcal S$. Jeśli $T$ jest ciągły,
to dla pewnych $j$ i $\varepsilon>0$ warunek
$Q_{j,j}(\varphi)<\varepsilon$ implikuje
$|T(\varphi)|<1$. Liniowość i skalowanie dają
$|T(\varphi)|\le (2/\varepsilon)Q_{j,j}(\varphi)$;
gdy seminorma znika, wniosek wynika ze skalowania przez
dowolną liczbę. Odwrotna implikacja jest natychmiastowa.
Włożenie $j:\mathcal D\hookrightarrow\mathcal S$ jest ciągłe:
dla niepustego zwartego $K\subseteq\R^d$ i
$\varphi\in\mathcal D_K$ mamy
$Q_{N,m}(\varphi)\le (1+\max_{x\in K}|x|)^N p_m(\varphi)$.
Dla $K=\varnothing$ etap jest przestrzenią zerową.
Ciągłość $j$ wynika z własności uniwersalnej granicy
induktywnej (Twierdzenie~\ref{thm:LF-universal18}).
Stąd ograniczenie $T|_{\mathcal D}=T\circ j$ jest ciągłe.
Odwzorowanie $j^{\mathrm t}:T\mapsto T\circ j$ jest
liniowe i różnowartościowe: jeśli dwa ciągłe funkcjonały
zgadzają się na gęstej podprzestrzeni $\mathcal D$,
to ich różnica znika na całej $\mathcal S$.
\end{proof}

\begin{twierdzenie}[Rodziny ograniczone dystrybucji temperowanych]
\label{thm:dt-tempered-bounded-families}
Niech $\mathcal A\subseteq\mathcal S'$. Następujące warunki
są równoważne:
\begin{enumerate}
\item $\mathcal A$ jest ograniczona w topologii słabej-*;
\item $\mathcal A$ jest ograniczona w topologii silnej;
\item istnieją $C>0$ i $N,m\in\N_0$ takie, że
\[
 |\langle T,\varphi\rangle|\le C Q_{N,m}(\varphi)
 \qquad(T\in\mathcal A,\ \varphi\in\mathcal S).
\]
\end{enumerate}
\end{twierdzenie}
\begin{proof}
Jeżeli $\mathcal A=\varnothing$, wszystkie warunki są
spełnione. W pozostałych przypadkach ograniczoność silna
implikuje słabą-* (singleton w $\mathcal S$ jest ograniczony).
Ograniczoność słaba-* oznacza, że dla każdej $\varphi\in\mathcal S$
zbiór $\{\langle T,\varphi\rangle:T\in\mathcal A\}$ jest
ograniczony w $\C$. Przestrzeń $\mathcal S$ jest przestrzenią
Frécheta, a zatem jest beczkowa
(Wniosek~\ref{cor:frechet-barrelled}).
Z twierdzenia Banacha--Steinhausa
(Twierdzenie~\ref{thm:banach-steinhaus-lcs}, z przestrzenią
wartości $\C$) wynika równociągłość rodziny $\mathcal A$. Istnieją więc
$j\in\N_0$ oraz $\varepsilon>0$ takie, że
$Q_{j,j}(\varphi)<\varepsilon$ implikuje
$|\langle T,\varphi\rangle|<1$ dla wszystkich $T\in\mathcal A$.
Dla $\varphi\ne0$ stosujemy to do
$\varepsilon\varphi/(2Q_{j,j}(\varphi))$.
Dostajemy wspólne oszacowanie z $C=2/\varepsilon$
oraz $N=m=j$; dla $\varphi=0$ nierówność jest oczywista.
Wreszcie wspólne oszacowanie implikuje
\[
 \sup_{T\in\mathcal A}P_B(T)
 \le C\sup_{\varphi\in B}Q_{N,m}(\varphi)<\infty
\]
dla każdego niepustego ograniczonego $B\subseteq\mathcal S$.
Jest to dokładnie ograniczoność w topologii silnej.
\end{proof}

\begin{uwaga}
Wspólne oszacowanie oznacza kontrolę zarówno wzrostu
w nieskończoności (parametr $N$), jak i liczby pochodnych
testu (parametr $m$). Nie można na ogół wybierać tych
parametrów osobno dla kolejnych elementów rodziny, jeżeli
chcemy wnioskować o jej ograniczoności. Przykładowo
rodzina $(\delta_{ne_1})_{n\ge1}$ spełnia wspólne
oszacowanie z $N=m=0$.
\end{uwaga}

\begin{przyklad}[Dystrybucje regularne i miary]
\label{ex:tempered-regular22}
Jeżeli $f\in L^1_{\mathrm{loc}}(\R^d)$ i
dla pewnych $C>0$, $M\ge0$ zachodzi
$|f(x)|\le C(1+|x|)^M$ prawie wszędzie,
to $T_f(\varphi)=\int_{\R^d}f(x)\varphi(x)\,\mathrm d\lambda_d(x)$
określa element $\mathcal S'$. Istotnie, dla całkowitego $N>M+d$ mamy
\[
 \int_{\R^d}|f(x)\varphi(x)|\,\mathrm d\lambda_d(x)
 \le C Q_{N,0}(\varphi)
 \int_{\R^d}(1+|x|)^{M-N}\,\mathrm d\lambda_d(x).
\]
Ostatnia całka jest skończona: na kuli jednostkowej
funkcja jest ograniczona, a na pierścieniu
$2^k\le |x|<2^{k+1}$ całkę ogranicza stała razy
$2^{k(d+M-N)}$. Szereg tych oszacowań jest zbieżny.
Z Twierdzenia~\ref{thm:tempered-estimate22} wynika
temperowalność $T_f$.
Bardziej ogólnie, jeżeli $\mu$ jest zespoloną lokalnie skończoną
miarą borelowską i dla pewnego $N\ge0$ zachodzi
$\int_{\R^d}(1+|x|)^{-N}\,\mathrm d|\mu|(x)<\infty$,
to $\langle T_\mu,\varphi\rangle=
\int_{\R^d}\varphi(x)\,\mathrm d\mu(x)$
definiuje dystrybucję temperowaną. Istotnie, biorąc
liczbę całkowitą $n\ge N$, mamy
\[
 \left|\int_{\R^d}\varphi\,\mathrm d\mu\right|
 \le Q_{n,0}(\varphi)
 \int_{\R^d}(1+|x|)^{-N}\,\mathrm d|\mu|(x).
\]
Całka jest bezwzględnie zbieżna względem $|\mu|$,
a oszacowanie dowodzi ciągłości funkcjonału.
Przykładami są funkcja stała $1$, wielomiany i miary Diraca.
\end{przyklad}

\begin{przyklad}[Dystrybucja nietemperowana]
\label{ex:non-tempered22}
Dystrybucja regularna wyznaczona przez
$f(x)=e^{|x|^2}$ należy do $\mathcal D'(\R^d)$,
ale nie jest ograniczeniem żadnego elementu $\mathcal S'$.
Aby to wykazać, weźmy $\chi\in\mathcal D(\R^d)$,
$0\le\chi\le1$, równą $1$ na $\overline B(0,1)$,
i funkcje $\varphi_R(x)=\chi(x/R)e^{-|x|^2/2}$.
Dla $R\to\infty$ zachodzi
$\varphi_R\to e^{-|x|^2/2}$ w $\mathcal S$
na mocy dowodu Twierdzenia~\ref{thm:density-schwartz17}:
dla $g(x)=e^{-|x|^2/2}\in\mathcal S$ i $R\ge1$ zachodzi
$Q_{N,m}(g-\chi(\cdot/R)g)\le
C_m(1+R)^{-1}Q_{N+1,m}(g)$.
Przynależność $g$ do $\mathcal S$ wynika stąd, że
każda jego pochodna jest wielomianem pomnożonym przez $g$.
Jednak
$\langle T_f,\varphi_R\rangle
=\int_{\R^d}\chi(x/R)e^{|x|^2/2}\,\mathrm d\lambda_d(x)
\to+\infty$,
gdyż dla $R\ge1$ całka jest nie mniejsza niż
$\int_{B(0,R)}1\,\mathrm d\lambda_d=
R^d\lambda_d(B(0,1))$.
Gdyby istniało ciągłe przedłużenie na $\mathcal S$,
wartości te musiałyby zbiegać do skończonej liczby.
\end{przyklad}

\begin{stwierdzenie}[Podprzestrzenie klasycznych dystrybucji]
\label{prop:schwartz-embeddings22}
Istnieją naturalne, różnowartościowe i ciągłe odwzorowania
\[
 \mathcal D\hookrightarrow\mathcal S\hookrightarrow
 L^2(\R^d)\hookrightarrow\mathcal S'\hookrightarrow\mathcal D',
\]
przy czym obie przestrzenie dualne wyposażamy albo
w topologie silne, albo w topologie słabe-*,
gdzie funkcję $f\in L^2$ identyfikujemy z funkcjonałem
$\langle T_f,\varphi\rangle=\int_{\R^d}f(x)\varphi(x)
\,\mathrm d\lambda_d(x)$.
Ponadto każda dystrybucja o zwartym nośniku jest temperowana:
$\mathcal E'(\R^d)\subseteq\mathcal S'(\R^d)$.
\end{stwierdzenie}
\begin{proof}
Pierwsze włożenie wykazano w Twierdzeniu~
\ref{thm:tempered-estimate22}. Drugie jest ciągłe na mocy
Stwierdzenia~\ref{prop:schwartz-operators17} i różnowartościowe,
ponieważ ciągła funkcja równa zeru prawie wszędzie znika
wszędzie. Dla $f\in L^2$ nierówność Cauchy'ego--Schwarza daje
$|\langle T_f,\varphi\rangle|\le\|f\|_2\|\varphi\|_2$;
także $\|\varphi\|_2\le C_dQ_{d+1,0}(\varphi)$.
Dla ograniczonego $B\subseteq\mathcal S$ otrzymujemy
\[
 P_B(T_f)\le C_d\|f\|_2\sup_{\varphi\in B}Q_{d+1,0}(\varphi).
\]
Prawa strona jest skończona i dowodzi ciągłości silnej,
a dla singletonów także słabej-*.
Jeśli $T_f=0$, to $\int f\varphi\,\mathrm d\lambda_d=0$
dla wszystkich $\varphi\in\mathcal D$.
Z gęstości $\mathcal D$ w $L^2$
(Twierdzenie~\ref{thm:dt-test-density}) wybieramy $\varphi_n\to
\overline f$ w $L^2$; nierówność Cauchy'ego--Schwarza daje
$\int |f|^2\,\mathrm d\lambda_d=0$, więc $f=0$ w $L^2$.
Ostatnie odwzorowanie jest różnowartościowe z
Twierdzenia~\ref{thm:tempered-estimate22}. Dla
$B\subseteq\mathcal D$ ograniczonego mamy
$P_B(T|_{\mathcal D})=P_{j(B)}(T)$, gdzie $j(B)$ jest
ograniczony w $\mathcal S$ dzięki ciągłości $j$.
Daje to ciągłość silną, a oceny na pojedynczych
funkcjach testowych dają ciągłość słabą-*.
Jeżeli $T\in\mathcal E'$ ma zwarty nośnik,
to jego przedłużenie z Lematu~\ref{lem:compact-action21}
spełnia $|T(\varphi)|\le C Q_{0,m}(\varphi)$.
Jest więc ciągłe na $\mathcal S$ i zgadza się z $T$
na $\mathcal D$; identyfikujemy $T$ z tym jednoznacznym
przedłużeniem.
\end{proof}

\subsection{Operatory na przestrzeni dystrybucji temperowanych}

\begin{twierdzenie}[Różniczkowanie i mnożenie przez wielomian]
\label{thm:tempered-operators22}
Dla każdego $\alpha,\beta\in\N_0^d$ operatory
$T\mapsto\partial^\alpha T$ i $T\mapsto x^\beta T$
są liniowe i ciągłe na $\mathcal S'$ zarówno dla topologii
$\sigma(\mathcal S',\mathcal S)$, jak i
$\beta(\mathcal S',\mathcal S)$, gdzie
\[
\langle\partial^\alpha T,\varphi\rangle
=(-1)^{|\alpha|}\langle T,\partial^\alpha\varphi\rangle,
\qquad\langle x^\beta T,\varphi\rangle
=\langle T,x^\beta\varphi\rangle.
\]
To samo dotyczy mnożenia przez każdą funkcję
$a\in C^\infty(\R^d;\C)$, dla której dla każdego
$\gamma\in\N_0^d$ istnieją $C_\gamma>0$ oraz
$M_\gamma\in\N_0$ takie, że
$|\partial^\gamma a(x)|\le C_\gamma(1+|x|)^{M_\gamma}$
dla wszystkich $x\in\R^d$.
\end{twierdzenie}
\begin{proof}
Różniczkowanie i mnożenie przez $x^\beta$ są ciągłymi
operatorami liniowymi $\mathcal S\to\mathcal S$,
co wykazano w Stwierdzeniu~\ref{prop:schwartz-operators17}.
Dla ogólnego mnożnika $a$ ustalmy $N,m\in\N_0$,
$M_m=\max_{|\gamma|\le m}M_\gamma$ i
$C_m=\max_{|\beta|\le m}\sum_{\gamma\le\beta}
\binom\beta\gamma C_\gamma$.
Ze wzoru Leibniza, dla $|\beta|\le m$, dostajemy
\[
 (1+|x|)^N|\partial^\beta(a\varphi)(x)|
 \le C_m Q_{N+M_m,m}(\varphi).
\]
Po wzięciu supremum oraz maksimum otrzymujemy
$Q_{N,m}(a\varphi)\le C_m Q_{N+M_m,m}(\varphi)$.
Zatem odwzorowanie $\varphi\mapsto a\varphi$ jest
ciągłym operatorem na $\mathcal S$.
Dla dowolnego ciągłego operatora liniowego
$A:\mathcal S\to\mathcal S$ jego transpozycja
$A^{\mathrm t}T=T\circ A$ spełnia
$P_B(A^{\mathrm t}T)=P_{A(B)}(T)$.
Ciągłość $A$ zapewnia ograniczoność $A(B)$, skąd
ciągłość silna; dla singletonów otrzymujemy ciągłość
słabą-*. Stosujemy to do $A\varphi=(-1)^{|\alpha|}
\partial^\alpha\varphi$ i $A\varphi=a\varphi$,
zgodnie z definicjami pochodnej i mnożenia dystrybucji.
Ograniczenie tych działań do $\mathcal D$ zgadza się
z działaniami na zwykłych dystrybucjach.
\end{proof}

\begin{twierdzenie}[Charakterystyka mnożników przestrzeni Schwartza]
\label{thm:dt-schwartz-multiplier-characterization}
Niech $a\in C^\infty(\R^d;\C)$. Następujące warunki
są równoważne:
\begin{enumerate}
\item $af\in\mathcal S$ dla każdego $f\in\mathcal S$;
\item operator $M_a:f\mapsto af$ jest ciągły na $\mathcal S$;
\item dla każdego $\alpha\in\N_0^d$ istnieją
$C_\alpha>0$ i $N_\alpha\in\N_0$ takie, że
\[
 |\partial^\alpha a(x)|\le C_\alpha(1+|x|)^{N_\alpha}
 \qquad(x\in\R^d).
\]
\end{enumerate}
\end{twierdzenie}
\begin{proof}
Implikację z warunku 3 do 2 wykazano w dowodzie
Twierdzenia~\ref{thm:tempered-operators22}, a warunek 2
implikuje 1 z definicji operatora na $\mathcal S$.
Pokażemy, że warunek 1 implikuje 3.
Gdyby warunek 3 nie zachodził, dla pewnego ustalonego
wielowskaźnika $\alpha$ i każdego $n\ge1$ funkcja
$x\mapsto (1+|x|)^{-n}|\partial^\alpha a(x)|$
byłaby nieograniczona. Ponieważ jest ciągła, jest
ograniczona na każdym zbiorze zwartym.
Dla $n=1$ wybieramy punkt poza kulą domkniętą
o promieniu $1$, gdzie iloraz przekracza $1$.
Dla $n\ge2$ wybieramy punkt poza kulą domkniętą
o promieniu $\max(n,|x_{n-1}|+2)$, gdzie odpowiedni
iloraz przekracza $1$. Jest to możliwe, ponieważ
ograniczoność ilorazu poza taką kulą wraz z jego
ciągłością na kuli dawałaby ograniczoność globalną.
Otrzymujemy zatem punkty $x_n$ takie, że
$|x_n|>n$, $|x_n|>|x_{n-1}|+2$ dla $n\ge2$ oraz
\[
 |\partial^\alpha a(x_n)|>(1+|x_n|)^n.
\]
Z Lematu~\ref{lem:cutoffs17} wybieramy
$\chi\in\mathcal D(B(0,1/2);\R)$ równą $1$ w otoczeniu zera
i przedłużamy ją zerem na $\R^d$.
Niech
\[
 f(x)=\sum_{n=1}^\infty
 (1+|x_n|)^{-n/2}\chi(x-x_n).
\]
Nośniki składników są parami rozłączne, ponieważ
$|x_n-x_k|\ge\bigl||x_n|-|x_k|\bigr|>2$ dla $n\ne k$.
Rodzina nośników jest lokalnie skończona dzięki
$|x_n|\to\infty$, zatem $f$ jest gładka.
Dla $|x-x_n|\le1/2$ mamy
$1+|x|\le(3/2)(1+|x_n|)$. W konsekwencji dla $N,m\in\N_0$
\[
 Q_{N,m}(f)\le (3/2)^N Q_{0,m}(\chi)
 \sup_{n\ge1}(1+|x_n|)^{N-n/2}<\infty.
\]
Istotnie, dla $n\ge2N$ wyrażenie pod supremum jest
nie większe niż $1$, a pozostałych indeksów jest
skończenie wiele. Poza sumą nośników wszystkie
pochodne $f$ znikają.
Stąd $f\in\mathcal S$. W otoczeniu $x_n$ funkcja $f$
jest stała i równa $(1+|x_n|)^{-n/2}$, więc
\[
 |\partial^\alpha(af)(x_n)|
 =(1+|x_n|)^{-n/2}|\partial^\alpha a(x_n)|
 >(1+|x_n|)^{n/2}\longrightarrow\infty.
\]
To przeczy $af\in\mathcal S$, gdyż $Q_{0,|\alpha|}(af)$
musiałaby być skończona. Otrzymujemy warunek 3.
\end{proof}

\begin{lemat}[Względna zwartość zbiorów ograniczonych w $\mathcal S$]
\label{lem:montel-schwartz22}
Każdy zbiór ograniczony w $\mathcal S(\R^d)$ ma zwarte
domknięcie w jej topologii Frécheta.
\end{lemat}
\begin{proof}
Niech $B\subseteq\mathcal S$ będzie ograniczony i
$(\varphi_n)\subseteq B$ dowolnym ciągiem.
Dla każdej liczby całkowitej $R\ge1$ i każdego $m\ge0$
rodziny pochodnych $\partial^\alpha\varphi_n$ dla
$|\alpha|\le m$ są jednostajnie ograniczone i
jednostajnie równociągłe na $\overline B(0,R)$:
równociągłość wynika z ograniczenia pochodnych
rzędu $m+1$ i całkowania wzdłuż odcinka w kuli:
dla $x,y$ w tej kuli oraz $|\alpha|\le m$ mamy
$|\partial^\alpha\varphi_n(x)-\partial^\alpha\varphi_n(y)|
\le\|x-y\|_1\sup_{\psi\in B}Q_{0,m+1}(\psi)$. Z twierdzenia Arzelà--Ascoliego
i argumentu diagonalnego uzyskujemy podciąg
$(\varphi_{n_k})$, którego wszystkie pochodne zbiegają
jednostajnie na wszystkich kulach domkniętych.
Nie musimy osobno identyfikować granic pochodnych:
poniższe oszacowanie da zbieżność w zupełnej
przestrzeni Schwartza.
Dla ustalonych $N,m$ i $R\ge1$ poza $B(0,R)$ zachodzi
\[
(1+|x|)^N|\partial^\alpha(\varphi_{n_k}-\varphi_{n_l})(x)|
\le R^{-1}\bigl(Q_{N+1,m}(\varphi_{n_k})
+Q_{N+1,m}(\varphi_{n_l})\bigr)\quad(|\alpha|\le m).
\]
Prawa strona jest jednostajnie mała przy dużym $R$;
wewnątrz kuli $B(0,R)$ zbieżność wszystkich pochodnych
jest jednostajna. Dlatego $(\varphi_{n_k})$ jest ciągiem
Cauchy'ego we wszystkich seminormach $Q_{N,m}$.
Zupełność $\mathcal S$ daje zbieżność do elementu
$\mathcal S$. Ponieważ przestrzeń jest metryzowalna,
ciągowa względna zwartość zbioru jest równoważna
zwartości jego domknięcia. Istotnie, dla ciągu $(g_n)$
w domknięciu $B$ wybieramy $f_n\in B$ w odległości
metrycznej mniejszej niż $1/n$ od $g_n$.
Zbieżny podciąg $(f_{n_k})$ daje zbieżny podciąg
$(g_{n_k})$ o tej samej granicy, należącej do domknięcia.
Zwartość wynika z kryterium ciągowego dla przestrzeni
metrycznych.
\end{proof}

\begin{twierdzenie}[Zbieżność słaba-* i silna ciągów]
\label{thm:seq-topologies-schwartz22}
Niech $(T_n)$ będzie ciągiem dystrybucji temperowanych,
a $T\in\mathcal S'$. Wtedy
\[
T_n\longrightarrow T\text{ w }\sigma(\mathcal S',\mathcal S)
\quad\Longleftrightarrow\quad
T_n\longrightarrow T\text{ w }\beta(\mathcal S',\mathcal S).
\]
Mimo to topologie $\sigma(\mathcal S',\mathcal S)$ i
$\beta(\mathcal S',\mathcal S)$ nie są równe.
\end{twierdzenie}
\begin{proof}
Implikacja ze zbieżności silnej jest natychmiastowa.
Załóżmy zbieżność słabą-*. Rodzina
$\{T_n-T:n\ge1\}$ jest punktowo ograniczona.
Z Twierdzenia~\ref{thm:dt-tempered-bounded-families}
otrzymujemy wspólne oszacowanie seminormowe tej rodziny,
a więc jej równociągłość.
Dla dowolnego niepustego ograniczonego $B\subseteq\mathcal S$
Lemat~\ref{lem:montel-schwartz22} daje zwartość
$K=\overline B$. Równociągłość funkcjonałów i
punktowa zbieżność do zera zapewniają zbieżność
jednostajną na $K$: dla danego $\varepsilon>0$
wybieramy otoczenie $V$ z
$\sup_n|\langle T_n-T,v\rangle|<\varepsilon/3$
dla $v\in V$. Możemy wybrać $V$ otwarte dzięki
wspólnemu oszacowaniu seminormowemu. Z pokrycia
$K$ zbiorami $\varphi+V$, $\varphi\in K$, wybieramy
skończone podpokrycie $\varphi_j+V$, $j=1,\ldots,r$.
Następnie dla dostatecznie dużego $n$ wszystkie
$|\langle T_n-T,\varphi_j\rangle|<\varepsilon/3$.
Dla każdego $\varphi\in B$ istnieje $j$ taki, że
$\varphi-\varphi_j\in V$; liniowość i nierówność
trójkąta dają
$|\langle T_n-T,\varphi\rangle|<2\varepsilon/3$.
Po wzięciu supremum otrzymujemy
$P_B(T_n-T)\le2\varepsilon/3<\varepsilon$.

Aby wykazać różność topologii, wybierzmy niezerową funkcję
$\psi\in\mathcal D$ o nośniku w $B(0,1)$ i zdefiniujmy
$\psi_n(x)=e^{-n^2}\psi(x-3ne_1)$ dla $n\ge1$.
Zbiór $B=\{\psi_n:n\ge1\}$ jest ograniczony w
$\mathcal S$: dla ustalonych $N,m$ mamy
$Q_{N,m}(\psi_n)\le e^{-n^2}(2+3n)^N Q_{0,m}(\psi)$,
co jest ograniczone niezależnie od $n$.
Jego liniowa otoczka jest nieskończenie wymiarowa,
ponieważ nośniki są parami rozłączne i każda $\psi_n$
jest niezerowa.
Otoczenie silne $\{T:P_B(T)<1/2\}$ nie zawiera
żadnego słabo-* otoczenia zera: dla dowolnie wielu
$\varphi_1,\ldots,\varphi_r\in\mathcal S$
można wybrać $\psi_n$ poza ich liniową otoczką.
Podprzestrzeń $M=\operatorname{span}\{\varphi_1,\ldots,
\varphi_r\}$ jest skończenie wymiarowa, a więc domknięta
w przestrzeni lokalnie wypukłej Hausdorffa $\mathcal S$.
Ponieważ $\psi_n\notin M$, twierdzenie o oddzielaniu
punktu od domkniętej podprzestrzeni
(Twierdzenie~\ref{thm:point-closed-subspace})
daje ciągły funkcjonał liniowy zerujący się na $M$
i niezerowy w $\psi_n$. Dzieląc go przez wartość
w $\psi_n$, otrzymujemy $T\in\mathcal S'$ taki, że
$T(\varphi_j)=0$ dla wszystkich $j$ oraz $T(\psi_n)=1$.
Należy on do dowolnego podstawowego słabo-* otoczenia
wyznaczonego przez te $\varphi_j$, ale nie do
$\{T:P_B(T)<1/2\}$. Ponieważ topologia silna jest
co najmniej tak mocna jak słaba-*, dowodzi to ich różności.
\end{proof}

\subsection{Rozszerzenie transformacji Fouriera}

\noindent\emph{Powiązania tematyczne w Bibliotece matematycznej:} \atlasref{PDE-08::def:pde8-distFT}{PDE-08: transformacja Fouriera dystrybucji}.


\begin{definicja}[Transformacja Fouriera dystrybucji temperowanej]
\label{def:fourier-tempered22}
Dla $T\in\mathcal S'(\R^d)$ definiujemy
$\widehat T=\mathcal F T\in\mathcal S'$ przez
\[
\langle\mathcal F T,\varphi\rangle
=\langle T,\mathcal F\varphi\rangle
\qquad(\varphi\in\mathcal S).
\]
Występuje tu transformacja Fouriera na $\mathcal S$
z Definicji~\ref{def:fourier-schwartz17}, z normalizacją
$(2\pi)^{-d/2}$ i jądrem $e^{-ix\cdot\xi}$.
Ponieważ $\mathcal F:\mathcal S\to\mathcal S$ jest ciągła
(Twierdzenie~\ref{thm:fourier-schwartz17}), złożenie
$T\circ\mathcal F$ jest ciągłym $\C$-liniowym funkcjonałem.
Definicja nie zawiera sprzężenia zespolonego ani
transformacji odwrotnej.
\end{definicja}

\begin{twierdzenie}[Transformacja Fouriera jest automorfizmem]
\label{thm:fourier-auto-tempered22}
Transformacja $\mathcal F:\mathcal S'\to\mathcal S'$
jest ciągłym izomorfizmem liniowo-topologicznym
zarówno dla topologii $\sigma(\mathcal S',\mathcal S)$,
jak i $\beta(\mathcal S',\mathcal S)$.
Jej odwrotność dana jest transpozycją
$\mathcal F^{-1}$ na $\mathcal S$:
$\langle\mathcal F^{-1}T,\varphi\rangle
=\langle T,\mathcal F^{-1}\varphi\rangle$.
Dla każdego $T\in\mathcal S'$ i $j\in\{1,\ldots,d\}$
zachodzą wzory
\[
\mathcal F(\partial_jT)=i\xi_j\,\mathcal F T,
\qquad
\mathcal F(x_jT)=i\partial_{\xi_j}(\mathcal F T),
\qquad\mathcal F^2T=T(-\cdot),
\]
gdzie $\langle T(-\cdot),\varphi\rangle=
\langle T,\varphi(-\cdot)\rangle$.
\end{twierdzenie}
\begin{proof}
Z Twierdzenia~\ref{thm:fourier-schwartz17} oba operatory
$\mathcal F$ i $\mathcal F^{-1}$ na $\mathcal S$ są
ciągłe. Dla $\varphi\in\mathcal S$ mamy
$\langle\mathcal F^{-1}\mathcal FT,\varphi\rangle
=\langle T,\mathcal F\mathcal F^{-1}\varphi\rangle
=\langle T,\varphi\rangle$.
Analogicznie
$\langle\mathcal F\mathcal F^{-1}T,\varphi\rangle
=\langle T,\mathcal F^{-1}\mathcal F\varphi\rangle
=\langle T,\varphi\rangle$.
Są to zatem wzajemnie odwrotne operatory na $\mathcal S'$.
Ciągłość słaba-* wynika z działania na pojedynczej
funkcji testowej. Dla silnej mamy
$P_B(\mathcal FT)=P_{\mathcal F(B)}(T)$;
$\mathcal F(B)$ jest ograniczony, gdy $B$ jest ograniczony,
co daje ciągłość. Odwrotność rozpatrujemy analogicznie.
Ustalmy $\varphi\in\mathcal S$, oznaczając jej argument
przez $\xi$, a argument $\mathcal F\varphi$ przez $x$.
Z Twierdzenia~\ref{thm:fourier-schwartz17} wynika
$\partial_{x_j}\mathcal F\varphi
=\mathcal F(-i\xi_j\varphi)$ oraz
$x_j\mathcal F\varphi=\mathcal F(-i\partial_{\xi_j}\varphi)$.
Zatem
\[
\begin{aligned}
\langle\mathcal F(\partial_jT),\varphi\rangle
&=-\langle T,\partial_{x_j}\mathcal F\varphi\rangle\\
&=\langle T,\mathcal F(i\xi_j\varphi)\rangle
=\langle i\xi_j\mathcal FT,\varphi\rangle,\\
\langle\mathcal F(x_jT),\varphi\rangle
&=\langle T,x_j\mathcal F\varphi\rangle\\
&=\langle T,\mathcal F(-i\partial_{\xi_j}\varphi)\rangle\\
&=-i\langle\mathcal FT,\partial_{\xi_j}\varphi\rangle
=\langle i\partial_{\xi_j}\mathcal FT,\varphi\rangle.
\end{aligned}
\]
Wzór odwrotny na $\mathcal S$ daje
$\mathcal F^{-1}\varphi(x)=\mathcal F\varphi(-x)$,
więc $\mathcal F^2\varphi(x)=\varphi(-x)$.
W konsekwencji
$\langle\mathcal F^2T,\varphi\rangle
=\langle T,\mathcal F^2\varphi\rangle
=\langle T,\varphi(-\cdot)\rangle$.
W szczególności $\mathcal F^4T=T$.
\end{proof}

\begin{przyklad}[Funkcja stała, delta Diraca i pochodne]
\label{ex:fourier-delta22}
Dla $a\in\R^d$ i $\alpha\in\N_0^d$ mamy w $\mathcal S'$
wzory
\[
\begin{aligned}
\mathcal F\delta_a&=(2\pi)^{-d/2}T_{e^{-ia\cdot\xi}},\\
\mathcal F T_1&=(2\pi)^{d/2}\delta_0,\\
\mathcal F(\partial^\alpha\delta_0)
&=(2\pi)^{-d/2}T_{(i\xi)^\alpha}.
\end{aligned}
\]
Przez $T_f$ oznaczamy dystrybucję regularną
reprezentowaną przez $f$, w szczególności $T_1$
odpowiada funkcji stałej $1$.
Pierwszy wynika z
$\langle\mathcal F\delta_a,\varphi\rangle
=(2\pi)^{-d/2}\int e^{-ia\cdot\xi}\varphi(\xi)
\,\mathrm d\lambda_d(\xi)$.
Dla $a=0$ mamy $\mathcal F\delta_0=(2\pi)^{-d/2}T_1$.
Ponieważ odbicie $\delta_0$ jest równe $\delta_0$,
zastosowanie $\mathcal F$ daje
$\delta_0=(2\pi)^{-d/2}\mathcal F T_1$.
Iterując wzór na transformatę pochodnej z
Twierdzenia~\ref{thm:fourier-auto-tempered22}, dostajemy
$\mathcal F(\partial^\alpha\delta_0)
=(i\xi)^\alpha\mathcal F\delta_0$, co daje trzeci wzór.
W szczególności delta Diraca nie jest funkcją,
ale jej transformata Fouriera jest funkcją gładką.
\end{przyklad}

\begin{stwierdzenie}[Zgodność z transformacją całkową]
\label{prop:dt-tempered-fourier-l1}
Jeżeli $f\in L^1(\R^d;\C)$, to zarówno $T_f$, jak i
$T_{\widehat f}$ są dystrybucjami temperowanymi, przy czym
\[
 \mathcal F T_f=T_{\widehat f}.
\]
Po prawej stronie $\widehat f$ jest transformatą całkową
z rozdziału o transformacji Fouriera funkcji całkowalnych.
\end{stwierdzenie}
\begin{proof}
Mamy $|\langle T_f,\varphi\rangle|
\le\|f\|_1 Q_{0,0}(\varphi)$, co daje temperowalność $T_f$.
Z Twierdzenia~\ref{thm:dt-fourier-l1-rl} funkcja $\widehat f$
jest ciągła i ograniczona, więc $T_{\widehat f}\in\mathcal S'$
na mocy Przykładu~\ref{ex:tempered-regular22}.
Dla $\varphi\in\mathcal S$ mamy, przy $c_d=(2\pi)^{-d/2}$,
\[
\begin{aligned}
 \langle\mathcal FT_f,\varphi\rangle
 &=\langle T_f,\mathcal F\varphi\rangle\\
 &=c_d\int_{\R^d}\int_{\R^d}f(x)e^{-ix\cdot\xi}
 \varphi(\xi)\,\mathrm d\lambda_d(\xi)\,\mathrm d\lambda_d(x)\\
 &=\int_{\R^d}\widehat f(\xi)\varphi(\xi)
 \,\mathrm d\lambda_d(\xi).
\end{aligned}
\]
Zamiana kolejności całkowania wynika z twierdzenia Fubiniego:
całka modułu z czynnikiem $c_d$ wynosi
$c_d\|f\|_1\|\varphi\|_1<\infty$.
Wobec tego obie definicje dają tę samą dystrybucję.
\end{proof}

\begin{stwierdzenie}[Przesunięcia i modulacje]
\label{prop:dt-tempered-fourier-shift-modulation}
Dla $h,\eta\in\R^d$ i $T\in\mathcal S'$ definiujemy
\[
 \langle\tau_h T,\varphi\rangle
 =\langle T,\varphi(\cdot+h)\rangle,\qquad
 M_\eta T=e^{i\eta\cdot x}T.
\]
Są to ciągłe automorfizmy liniowe w obu topologiach
przestrzeni $\mathcal S'$. Zachodzą tożsamości
\[
 \mathcal F(\tau_h T)=e^{-ih\cdot\xi}\mathcal FT,
 \qquad \mathcal F(M_\eta T)=\tau_\eta(\mathcal FT).
\]
\end{stwierdzenie}
\begin{proof}
Dla przesunięcia testu mamy
$Q_{N,m}(\varphi(\cdot+h))\le
(1+|h|)^NQ_{N,m}(\varphi)$, ponieważ
$1+|x|\le(1+|h|)(1+|x+h|)$.
Ciągłość mnożenia przez $e^{i\eta\cdot x}$ wynika z
Twierdzenia~\ref{thm:tempered-operators22}, gdyż wszystkie
pochodne tego mnożnika są ograniczone. Ciągłość transpozycji
wynika z tego samego twierdzenia. Operatory odwrotne to
odpowiednio $\tau_{-h}$ oraz $M_{-\eta}$.
Dla dystrybucji regularnej przesunięcie odpowiada
funkcji $x\mapsto f(x-h)$, przez podstawienie w całce.
Z definicji transformacji testów bezpośrednio otrzymujemy
\[
 (\mathcal F\varphi)(x+h)
 =\mathcal F(e^{-ih\cdot\xi}\varphi)(x),
 \qquad
 e^{i\eta\cdot x}(\mathcal F\varphi)(x)
 =\mathcal F(\varphi(\cdot+\eta))(x).
\]
Obie całki są bezwzględnie zbieżne, a druga równość
wynika z podstawienia $\zeta=\xi+\eta$.
Stosując te równości w definicji transpozycji, dostajemy
\[
\begin{aligned}
 \langle\mathcal F(\tau_hT),\varphi\rangle
 &=\langle T,(\mathcal F\varphi)(\cdot+h)\rangle
 =\langle\mathcal FT,e^{-ih\cdot\xi}\varphi\rangle,\\
 \langle\mathcal F(M_\eta T),\varphi\rangle
 &=\langle T,e^{i\eta\cdot x}\mathcal F\varphi\rangle
 =\langle\mathcal FT,\varphi(\cdot+\eta)\rangle .
\end{aligned}
\]
To dowodzi obu wzorów.
\end{proof}

\begin{przyklad}[Fale płaskie]
Dla $a,\eta\in\R^d$ mamy
\[
 \mathcal F\delta_a=c_d T_{e^{-ia\cdot\xi}},
 \qquad
 \mathcal F T_{e^{i\eta\cdot x}}=c_d^{-1}\delta_\eta,
 \qquad c_d=(2\pi)^{-d/2}.
\]
Pierwszy wzór wynika z obliczenia
$\langle\mathcal F\delta_a,\varphi\rangle
=(\mathcal F\varphi)(a)$.
W drugim używamy $\mathcal F T_1=c_d^{-1}\delta_0$
z Przykładu~\ref{ex:fourier-delta22} i poprzedniego
stwierdzenia. Fala o jednej częstotliwości ma więc
transformatę skupioną w jednym punkcie, mimo że nie
jest funkcją całkowalną na $\R^d$.
\end{przyklad}

\begin{stwierdzenie}[Odwracalna liniowa zmiana współrzędnych]
\label{prop:dt-tempered-linear-change}
Niech $A\in GL_d(\R)$ i $T\in\mathcal S'(\R^d)$.
Definiujemy $L_A T$ przez
\[
 \langle L_A T,\varphi\rangle
 =|\det A|^{-1}\langle T,\varphi\circ A^{-1}\rangle,
 \qquad \varphi\in\mathcal S.
\]
Operator $L_A$ jest automorfizmem liniowo-topologicznym
w obu topologiach $\mathcal S'$, a jego odwrotnością
jest $L_{A^{-1}}$. Jeżeli $f\in L^1_{\mathrm{loc}}(\R^d;\C)$ i dystrybucja
regularna $T_f$ jest temperowana, to $T_{f\circ A}$
jest temperowana oraz $L_A T_f=T_{f\circ A}$.
Obie dystrybucje regularne rozumiemy tu jako
jednoznaczne ciągłe przedłużenia z $\mathcal D$ na $\mathcal S$.
Zachodzi
\[
 \mathcal F(L_A T)=|\det A|^{-1}L_{A^{-T}}(\mathcal F T).
\]
\end{stwierdzenie}
\begin{proof}
Dla każdej macierzy odwracalnej $B$ operator
$C_B\varphi=\varphi\circ B$ zachowuje $\mathcal S$
i jest ciągły. Istotnie, z reguły łańcuchowej każda
pochodna rzędu $k$ funkcji $\varphi\circ B$ jest
skończoną kombinacją pochodnych rzędu $k$ funkcji
$\varphi$, obliczanych w punkcie $Bx$, ze stałymi
współczynnikami zależnymi od $B$.
Ponadto $1+|x|\le\max(1,\|B^{-1}\|)(1+|Bx|)$.
Stąd dla $N,m\in\N_0$ istnieje stała $C_{B,N,m}$ taka,
że $Q_{N,m}(C_B\varphi)\le C_{B,N,m}Q_{N,m}(\varphi)$.
Zastosowanie tego argumentu do $B^{-1}$ daje ciągłość
odwrotności. Definicja $L_A$ jest więc transpozycją
ciągłego operatora z czynnikiem skalarnym; ciągłość
w obu topologiach wynika z argumentu w dowodzie
Twierdzenia~\ref{thm:tempered-operators22}.
Bezpośrednie złożenie definicji dla $A$ i $A^{-1}$
daje operator identycznościowy w obu kolejnościach.
Dla $f\in L^1_{\mathrm{loc}}$ podstawienie $y=Ax$
na zbiorach zwartych pokazuje lokalną całkowalność
$f\circ A$, a dla $\varphi\in\mathcal D$ daje
\[
 \langle L_A T_f,\varphi\rangle
 =|\det A|^{-1}\int_{\R^d}f(y)\varphi(A^{-1}y)
 \,\mathrm d\lambda_d(y)
 =\int_{\R^d}f(Ax)\varphi(x)\,\mathrm d\lambda_d(x).
\]
Zatem $L_A T_f$ jest ciągłym przedłużeniem dystrybucji
regularnej $T_{f\circ A}$ na $\mathcal S$.
Jednoznaczność wynika z Twierdzenia~\ref{thm:tempered-estimate22}.
Nie twierdzimy, że całka z $f\varphi$ jest bezwzględnie
zbieżna dla każdego $\varphi\in\mathcal S$, gdy jedynym
założeniem o wzroście $f$ jest temperowalność $T_f$.

Z reguły zmiany zmiennych w transformacji całkowej
(Stwierdzenie~\ref{prop:dt-fourier-l1-rules}) wynika
\[
 \mathcal F(\varphi\circ A^T)(x)
 =|\det A|^{-1}(\mathcal F\varphi)(A^{-1}x).
\]
Wobec tego
\[
\begin{aligned}
 \langle\mathcal F(L_AT),\varphi\rangle
 &=|\det A|^{-1}
 \langle T,(\mathcal F\varphi)\circ A^{-1}\rangle\\
 &=\langle T,\mathcal F(\varphi\circ A^T)\rangle\\
 &=\langle\mathcal F T,\varphi\circ A^T\rangle\\
 &=|\det A|^{-1}\langle L_{A^{-T}}(\mathcal FT),\varphi\rangle .
\end{aligned}
\]
Ostatni krok wykorzystuje
$|\det A^{-T}|^{-1}=|\det A|$. Dowodzi to tezy.
\end{proof}

\begin{przyklad}[Transformata wielomianu]
\label{ex:dt-tempered-polynomial}
Dla $\alpha\in\N_0^d$ zachodzi
\[
 \mathcal F T_{x^\alpha}
 =(2\pi)^{d/2}i^{|\alpha|}\partial^\alpha\delta_0.
\]
Wielomian wyznacza dystrybucję temperowaną na mocy
Przykładu~\ref{ex:tempered-regular22}.
Z Przykładu~\ref{ex:fourier-delta22} mamy
$\mathcal FT_1=(2\pi)^{d/2}\delta_0$.
Ponieważ $x^\alpha T_1=T_{x^\alpha}$, iteracja
wzoru na transformatę mnożenia przez współrzędną
z Twierdzenia~\ref{thm:fourier-auto-tempered22}
daje podaną równość.
\end{przyklad}

\subsection{Splot z funkcją szybko malejącą}

\begin{twierdzenie}[Splot z elementem przestrzeni Schwartza]
\label{thm:schwartz-convolution22}
Niech $T\in\mathcal S'(\R^d)$ i
$\rho\in\mathcal S(\R^d)$. Funkcja
$u(x)=\langle T_y,\rho(x-y)\rangle$ należy do
$C^\infty(\R^d)$, a każda jej pochodna ma
co najwyżej wzrost wielomianowy. Dystrybucja
regularna $T_u$ jest temperowana i oznaczamy ją $T*\rho$.
Dla każdego $\alpha\in\N_0^d$ zachodzi
\[
\mathcal F(T*\rho)=(2\pi)^{d/2}(\mathcal FT)\,(\mathcal F\rho),
\qquad \partial^\alpha(T*\rho)=T*(\partial^\alpha\rho).
\]
\end{twierdzenie}
\begin{proof}
Dla $h\in\R^d$ oraz $N,m\in\N_0$ zachodzi
$Q_{N,m}(\rho(h-\cdot))\le(1+|h|)^N Q_{N,m}(\rho)$,
gdyż $1+|y|\le(1+|h|)(1+|h-y|)$.
Ta sama nierówność obowiązuje dla każdej pochodnej
$\rho$. Z twierdzenia Newtona--Leibniza zastosowanego
dwukrotnie, dla $0<|t|\le1$ i $j\in\{1,\ldots,d\}$,
otrzymujemy
\[
\begin{aligned}
 Q_{N,m}\left(
 \frac{\rho(x+te_j-\cdot)-\rho(x-\cdot)}{t}
 -(\partial_j\rho)(x-\cdot)\right)
 \le |t|(2+|x|)^N Q_{N,m+2}(\rho).
\end{aligned}
\]
Istotnie, po każdej pochodnej względem $y$ resztę
ograniczamy przez $|t|$ razy supremum odpowiedniej
pochodnej drugiego rzędu względem parametru na
odcinku od $x$ do $x+te_j$ i stosujemy poprzednią
nierówność wagową. Jednokrotne zastosowanie twierdzenia
Newtona--Leibniza daje również ciągłość przesunięć
w każdej seminormie. Powtarzając argument dla
pochodnych $\rho$, otrzymujemy gładkość funkcji
$x\mapsto\rho(x-\cdot)$ o wartościach w $\mathcal S$.
Ciągłość $T$ pozwala więc różniczkować pod parowaniem:
\[
 \partial^\alpha u(x)=\langle T_y,
 (\partial^\alpha\rho)(x-y)\rangle.
\]
Z Twierdzenia~\ref{thm:tempered-estimate22} wybieramy
$C>0$, $N,m\in\N_0$ z $|T(\varphi)|\le C Q_{N,m}(\varphi)$.
Poprzednie oszacowanie przesunięcia daje
\[
 |\partial^\alpha u(x)|\le
 C(1+|x|)^N Q_{N,m+|\alpha|}(\rho).
\]
W szczególności $T_u\in\mathcal S'$ na mocy
Przykładu~\ref{ex:tempered-regular22}.

Ustalmy $\psi\in\mathcal S$ i niech $G_x(y)=
\psi(x)\rho(x-y)$. Odwzorowanie $x\mapsto G_x$
jest ciągłe w $\mathcal S$ i
\[
 Q_{N,m}(G_x)\le |\psi(x)|(1+|x|)^N Q_{N,m}(\rho).
\]
Prawa strona jest całkowalna po $\R^d$ dla każdych $N,m$.
Dla $R\ge1$ określmy $I_R(y)=\int_{[-R,R]^d}G_x(y)
\,\mathrm d\lambda_d(x)$. Różniczkowanie pod całką na
tym zwartym prostokącie jest dozwolone, a powyższe
oszacowanie dowodzi, że $I_R\in\mathcal S$.
Sumy Riemanna tej całki zbiegają do $I_R$ w każdej
seminormie $Q_{N,m}$: błąd jest nie większy niż
objętość prostokąta razy supremum
$Q_{N,m}(G_x-G_{x'})$ dla $x,x'$ z tej samej części
podziału. Supremum dąży do zera z jednostajnej
ciągłości w danej seminormie na prostokącie.
Stosując $T$ do sum Riemanna i przechodząc do granicy
otrzymujemy
\[
 \langle T,I_R\rangle=
 \int_{[-R,R]^d}\psi(x)u(x)\,\mathrm d\lambda_d(x).
\]
Dla $R'<R''$ oszacowanie ogona ma postać
\[
 Q_{N,m}(I_{R''}-I_{R'})
 \le Q_{N,m}(\rho)
 \int_{\R^d\setminus[-R',R']^d}
 (1+|x|)^N|\psi(x)|\,\mathrm d\lambda_d(x).
\]
W szczególności $I_n$ jest ciągiem Cauchy'ego w $\mathcal S$.
Z jej zupełności granica istnieje w $\mathcal S$, a
zbieżność punktowa identyfikuje ją z funkcją
$\check\rho*\psi$, gdzie $\check\rho(y)=\rho(-y)$:
\[
 (\check\rho*\psi)(y)=
 \int_{\R^d}\rho(x-y)\psi(x)\,\mathrm d\lambda_d(x).
\]
Z ciągłości $T$ oraz bezwzględnej całkowalności
$u\psi$ otrzymujemy
$\langle T_u,\psi\rangle=\langle T,\check\rho*\psi\rangle$.
Stwierdzenie~\ref{prop:convolution-schwartz17} zapewnia
ciągłość operatora $\psi\mapsto\check\rho*\psi$ na
$\mathcal S$. Jego transpozycją jest więc ciągły
operator $T\mapsto T*\rho$ dla obu topologii na
$\mathcal S'$, zgodnie z dowodem
Twierdzenia~\ref{thm:tempered-operators22}.
Jeśli $\rho\in\mathcal D$, konstrukcja zgadza się na
$\mathcal D$ z Twierdzeniem~\ref{thm:smoothing21}.

Pozostaje obliczyć stałą i znak w formule Fouriera.
Dla $\psi\in\mathcal S$ mamy, przy $c_d=(2\pi)^{-d/2}$,
\[
\begin{aligned}
(\check\rho*\mathcal F\psi)(y)
 &=\int_{\R^d}\rho(z)(\mathcal F\psi)(y+z)
 \,\mathrm d\lambda_d(z)\\
 &=c_d\int_{\R^d}e^{-iy\cdot\xi}\psi(\xi)
 \int_{\R^d}\rho(z)e^{-iz\cdot\xi}\,\mathrm d\lambda_d(z)
 \,\mathrm d\lambda_d(\xi)\\
 &=(2\pi)^{d/2}\mathcal F((\mathcal F\rho)\psi)(y).
\end{aligned}
\]
Twierdzenie Fubiniego stosujemy, ponieważ całka
modułu podwójnego jest ograniczona przez
$c_d\|\rho\|_1\|\psi\|_1$. Iloczyn
$(\mathcal F\rho)\psi$ należy do $\mathcal S$.
Zatem
\[
\begin{aligned}
\langle\mathcal F(T*\rho),\psi\rangle
&=\langle T,\check\rho*\mathcal F\psi\rangle\\
&=(2\pi)^{d/2}\langle\mathcal FT,
 (\mathcal F\rho)\psi\rangle.
\end{aligned}
\]
To jest pierwszy wzór, a drugi wynika ze wzoru na
$\partial^\alpha u$ i zgodności pochodnej dystrybucji
regularnej z klasyczną pochodną gładkiej funkcji
(wystarczy sprawdzić na $\mathcal D$, gęstej w $\mathcal S$).
\end{proof}

\begin{twierdzenie}[Aproksymacja dystrybucji temperowanych funkcjami testowymi]
\label{thm:dt-tempered-test-density}
Dla każdej $T\in\mathcal S'(\R^d)$ istnieje ciąg
$(f_n)\subseteq\mathcal D(\R^d)$ taki, że
$T_{f_n}\to T$ w silnej topologii $\beta(\mathcal S',\mathcal S)$.
Co więcej, ciąg można wybrać tak, by dla każdego
niepustego ograniczonego $B\subseteq\mathcal S$
zachodziło $P_B(T_{f_n}-T)=O(n^{-1})$.
\end{twierdzenie}
\begin{proof}
Niech $\chi\in\mathcal D(\R^d;\R)$ będzie równa $1$
na $\overline B(0,1)$, $0\le\chi\le1$,
$\supp\chi\subseteq B(0,2)$, jak w dowodzie
Twierdzenia~\ref{thm:density-schwartz17}.
Piszemy $\chi_n(x)=\chi(x/n)$. Wybieramy standardową
aproksymację jedności $\rho_{1/n}$ z
Definicji~\ref{def:mollifier21}.
Dystrybucja $\chi_nT$ ma zwarty nośnik, ponieważ jej
ograniczenie do $\mathcal D$ znika poza $\supp\chi_n$
na mocy definicji mnożenia przez funkcję gładką.
Z Twierdzenia~\ref{thm:smoothing21} splot
\[
 T_{f_n}=(\chi_nT)*\rho_{1/n}
\]
jest reprezentowany przez funkcję testową $f_n$,
gdyż jego nośnik zawiera się w zwartym zbiorze
$\supp\chi_n+\supp\rho_{1/n}$.
Z dowodu Twierdzenia~\ref{thm:schwartz-convolution22},
dla $\varphi\in\mathcal S$, mamy
\[
 \langle T_{f_n}-T,\varphi\rangle
 =\langle T,\chi_n(\check\rho_{1/n}*\varphi)-\varphi\rangle.
\]
Niech $N,m\in\N_0$. Dla $0<\varepsilon\le1$,
z normalizacji $\rho$ i podstawienia zmiennej otrzymujemy
\[
 (\check\rho_\varepsilon*\varphi)(x)-\varphi(x)
 =\int_{\R^d}\rho(z)
 [\varphi(x+\varepsilon z)-\varphi(x)]\,\mathrm d\lambda_d(z).
\]
Wzór pozostaje prawdziwy po zastosowaniu każdej pochodnej
$\partial^\alpha$ do obu stron, ponieważ całkujemy po
zwartym nośniku $\rho$, a pochodne $\varphi$ są ograniczone.
Dla $|z|\le1$, $0\le t\le1$ zachodzi
$1+|x|\le2(1+|x+t\varepsilon z|)$.
Stosując twierdzenie Newtona--Leibniza do różnicy
$\partial^\alpha\varphi(x+\varepsilon z)
-\partial^\alpha\varphi(x)$, dla $|\alpha|\le m$ dostajemy
\[
 Q_{N,m}(\check\rho_\varepsilon*\varphi-\varphi)
 \le 2^N\varepsilon Q_{N,m+1}(\varphi)
 \int_{\R^d}\|z\|_1\rho(z)\,\mathrm d\lambda_d(z),
\]
gdzie $\|z\|_1=\sum_{j=1}^d|z_j|$.
Z reguły Leibniza dla $|\alpha|\le m$ mamy
\[
 (1+|x|)^N|\partial^\alpha(\chi_n g)(x)|
 \le Q_{N,m}(g)\sum_{\gamma\le\alpha}
 \binom{\alpha}{\gamma}n^{-|\gamma|}
 \|\partial^\gamma\chi\|_\infty.
\]
Ponieważ $n^{-|\gamma|}\le1$ dla $n\ge1$, po wzięciu
supremum i maksimum otrzymujemy
$Q_{N,m}(\chi_n g)\le C_m Q_{N,m}(g)$
dla $g\in\mathcal S$, gdzie $C_m$ nie zależy od $n$. Z dowodu
Twierdzenia~\ref{thm:density-schwartz17} mamy też
$Q_{N,m}(\chi_n\varphi-\varphi)
\le C_m'(1+n)^{-1}Q_{N+1,m}(\varphi)$.
Rozkładając różnicę w parowaniu na
$\chi_n(\check\rho_{1/n}*\varphi-\varphi)
+(\chi_n\varphi-\varphi)$, otrzymujemy zatem
\[
\begin{aligned}
 &Q_{N,m}(\chi_n(\check\rho_{1/n}*\varphi)-\varphi)\\
 &\qquad\le \frac{C_{N,m,\rho,\chi}}{n}
 \bigl(Q_{N,m+1}(\varphi)+Q_{N+1,m}(\varphi)\bigr).
\end{aligned}
\]
Z Twierdzenia~\ref{thm:tempered-estimate22} wybieramy
$N,m,C$ takie, że $|T(g)|\le C Q_{N,m}(g)$.
Dla dowolnego niepustego zbioru ograniczonego
$B\subseteq\mathcal S$ obie seminormy po prawej stronie
ostatniego oszacowania są jednostajnie ograniczone na $B$.
Dokładniej, dla każdego takiego $B$ istnieje stała
$C_{T,B,\rho,\chi}\ge0$ niezależna od $n$, taka że
\[
 P_B(T_{f_n}-T)\le \frac{C_{T,B,\rho,\chi}}{n}
 \qquad(n\ge1).
\]
W szczególności $P_B(T_{f_n}-T)\to0$.
Jest to dokładnie zbieżność silna.
\end{proof}

\subsection{Łańcuch przestrzeni sprzężonych i kryterium temperowania}

\begin{lemat}[Transpozycja gęstego zanurzenia]
\label{lem:transpozycja25}
Niech $X,Y$ będą przestrzeniami lokalnie wypukłymi
Hausdorffa, a $j:X\to Y$ ciągłym odwzorowaniem liniowym
o gęstym obrazie. Wówczas transpozycja
$j':Y'\to X'$, $j'(T)=T\circ j$, jest liniowa,
różnowartościowa i ciągła zarówno dla topologii słabych-*, jak i dla
silnych topologii przestrzeni dualnych.
\end{lemat}
\begin{proof}
Dla każdego $T\in Y'$ złożenie $T\circ j$ jest ciągłym
funkcjonałem liniowym na $X$, więc $j'$ jest dobrze
określone; liniowość wynika z liniowości składania
względem pierwszego czynnika. Jeżeli $T\circ j=0$,
to $T=0$ przez gęstość $j(X)$ i ciągłość $T$.
Dla $x\in X$ mamy
$|\langle j'T,x\rangle|=|\langle T,jx\rangle|$,
co daje ciągłość w topologiach słabych-*.
Jeśli $B\subseteq X$ jest ograniczony, to $j(B)$ jest
ograniczony w $Y$ (dla każdej ciągłej seminormy $p$
na $Y$ seminorma $p\circ j$ jest ciągła na $X$).
Dla seminorm silnych zachodzi dokładna równość
$\sup_{x\in B}|\langle j'T,x\rangle|
= \sup_{y\in j(B)}|\langle T,y\rangle|$.
To dowodzi ciągłości w topologiach silnych.
\end{proof}

\begin{wniosek}[Łańcuch przestrzeni sprzężonych]
\label{cor:dual-lancuch25}
Dla każdego $d\ge1$ zachodzą kanoniczne, ciągłe i
różnowartościowe odwzorowania liniowe
\[
 \mathcal E'(\R^d)\longrightarrow\mathcal S'(\R^d)
 \longrightarrow\mathcal D'(\R^d),
\]
zarówno w topologiach słabych-*, jak i silnych.
Dla każdego niepustego zbioru otwartego $\Omega\subseteq\R^d$
kanoniczne odwzorowanie
$\mathcal E'(\Omega)\to\mathcal D'(\Omega)$
ma te same własności. Samo stwierdzenie ciągłości i
różnowartościowości \emph{nie oznacza}, że są to zanurzenia
topologiczne na obrazy z topologią podprzestrzeni.
\end{wniosek}
\begin{proof}
Stosujemy Lemat~\ref{lem:transpozycja25} do odwzorowań
z Twierdzenia~\ref{thm:zanurzenia25}. Jeżeli $i:X\to Y$
oraz $j:Y\to Z$ są tymi kolejnymi odwzorowaniami,
to $(j\circ i)'=i'\circ j'$, gdyż dla $T\in Z'$
mamy $(T\circ j)\circ i=T\circ(j\circ i)$.
Dlatego ograniczenie z $\mathcal E'$ do $\mathcal D'$
przez $\mathcal S'$ jest tym samym odwzorowaniem co
bezpośrednie ograniczenie. W przypadku zespolonym
transpozycja jest liniowa, bez sprzężenia zespolonego.
Ciągłość i różnowartościowość nie zapewniają ciągłości
odwrotności na obrazie; poniżej podamy kontrprzykłady.
\end{proof}
\subsection{Dystrybucje temperowane jako podprzestrzeń \texorpdfstring{$\mathcal D'$}{D'}}

\begin{twierdzenie}[Kryterium temperowania]
\label{thm:temp-criterion25}
Niech $d\ge1$ będzie liczbą całkowitą, $\K\in\{\R,\C\}$
i $T\in\mathcal D'(\R^d)$ nad $\K$.
Następujące warunki są równoważne:
\begin{enumerate}[label=\textup{(\roman*)}]
\item $T$ należy do obrazu kanonicznego odwzorowania
$\mathcal S'(\R^d)\to\mathcal D'(\R^d)$;
\item istnieją $C>0$, $N,m\in\N_0$ takie, że
$|\langle T,\varphi\rangle|\le C p_{N,m}(\varphi)$
dla każdej $\varphi\in\mathcal D(\R^d)$.
\end{enumerate}
Jeżeli te warunki zachodzą, przedłużenie
$\widetilde T\in\mathcal S'(\R^d)$ jest jednoznaczne
i spełnia $|\langle\widetilde T,f\rangle|\le C p_{N,m}(f)$
dla wszystkich $f\in\mathcal S(\R^d)$, z tymi samymi
stałymi co w~\textup{(ii)}. W szczególności każda
dystrybucja o zwartym nośniku jest temperowana.
\end{twierdzenie}
\begin{proof}
Niech najpierw $\widetilde T\in\mathcal S'$ będzie
przedłużeniem $T$. Z ciągłości w zerze i współkońcowości
rosnącego ciągu norm $p_{j,j}$ istnieją $j\in\N_0$ oraz
$\delta>0$ takie, że $p_{j,j}(f)<\delta$ implikuje
$|\widetilde T(f)|<1$. Dla $f\ne0$ stosujemy ten warunek
do $\delta f/(2p_{j,j}(f))$, otrzymując
$|\widetilde T(f)|\le (2/\delta)p_{j,j}(f)$.
Dla $f=0$ nierówność jest oczywista. Ograniczenie
do $\mathcal D$ daje warunek~\textup{(ii)}.
Jest to argument z Twierdzenia~\ref{thm:tempered-estimate22},
działający nad obydwoma ciałami.

Załóżmy teraz~\textup{(ii)}. Dla każdego $f\in\mathcal S$
wybieramy ciąg $(\varphi_j)$ w $\mathcal D$ zbieżny do $f$
w $\mathcal S$ (Twierdzenie~\ref{thm:zanurzenia25}).
Ponieważ
\[
 |T(\varphi_j)-T(\varphi_k)|
 \le C p_{N,m}(\varphi_j-\varphi_k)\longrightarrow0
 \quad(j,k\to\infty),
\]
ciąg wartości ma granicę w zupełnym ciele $\K$.
Kładziemy $\widetilde T(f)=\lim_jT(\varphi_j)$.
Dla drugiego ciągu $(\psi_j)$ zbieżnego do $f$ mamy
$|T(\varphi_j)-T(\psi_j)|\le
C p_{N,m}(\varphi_j-\psi_j)\to0$, więc definicja nie zależy
od wyboru ciągu. Dla $f,g\in\mathcal S$ i $a,b\in\K$
można aproksymować $af+bg$ przez $a\varphi_j+b\psi_j$,
gdzie teraz $\varphi_j\to f$ i $\psi_j\to g$.
Przechodząc do granicy, otrzymujemy liniowość.
Ponadto $|p_{N,m}(\varphi_j)-p_{N,m}(f)|
\le p_{N,m}(\varphi_j-f)$, więc
\[
 |\widetilde T(f)|\le C p_{N,m}(f).
\]
Daje to ciągłość; na $\mathcal D$ wystarczy wziąć
ciąg stały, aby uzyskać $\widetilde T|_{\mathcal D}=T$.
Jednoznaczność wynika z gęstości $\mathcal D$.

Jeżeli $T$ ma zwarty nośnik, to z
Twierdzenia~\ref{thm:compact-support-E25} otrzymujemy
$|T(\varphi)|\le C p_{K,m}(\varphi)\le C p_{0,m}(\varphi)$
na funkcjach testowych, a zatem~\textup{(ii)} zachodzi z $N=0$.
Przedłużenie jest ograniczeniem funkcjonału na $\mathcal E$
do $\mathcal S$, przez jednoznaczność; odpowiada to
Stwierdzeniu~\ref{prop:schwartz-embeddings22}.
\end{proof}

\begin{przyklad}[Dystrybucja nietemperowana]
\label{ex:non-tempered25}
Na $\R$ funkcja $f(x)=e^x$ jest lokalnie całkowalna,
więc wyznacza dystrybucję regularną $T_f\in\mathcal D'(\R)$.
Nie jest ona temperowana. Wybierzmy nieujemną, niezerową
$\psi\in\mathcal D((-1,1))$, przedłużoną zerem na $\R$,
i niech $\varphi_n(x)=e^{-n/2}\psi(x-n)$ dla $n\ge1$.
Dla $x\in\supp\varphi_n$ mamy $|x|\le n+1$; pochodne
$\varphi_n$ znikają poza tym nośnikiem. Stąd dla wszystkich
$N,m\in\N_0$
\[
 p_{N,m}(\varphi_n)\le e^{-n/2}(n+2)^N p_{0,m}(\psi)
 \longrightarrow0,
\]
gdyż funkcja wykładnicza rośnie szybciej niż każdy
wielomian. Jednak po podstawieniu $y=x-n$ otrzymujemy
\[
 \langle T_f,\varphi_n\rangle=
 e^{n/2}\int_\R e^y\psi(y)\,\mathrm d\lambda_1(y)
 \longrightarrow+\infty.
\]
Całka jest skończona i ściśle dodatnia: $\psi$ jest
nieujemna i dodatnia na pewnym niepustym przedziale.
Ciągłość przedłużenia na $\mathcal S$ wymagałaby
$\langle T_f,\varphi_n\rangle\to0$, co daje sprzeczność.
Przykład działa nad $\R$ i $\C$; razem z
Przykładem~\ref{ex:non-tempered22} pokazuje on dwie
metody wykrywania braku temperowalności (przesuwanie
nośnika oraz powiększanie odcięcia).
\end{przyklad}
\begin{przyklad}[Zbieżność i ucieczka nośnika]
\label{ex:escaping-deltas25}
Na $\R$ niech $\delta_n$ oznacza deltę Diraca w punkcie
$n$ dla $n\ge1$. Każda z tych dystrybucji ma zwarty
nośnik, więc należy także do $\mathcal E'$ i $\mathcal S'$
na mocy Twierdzeń~\ref{thm:compact-support-E25}
i~\ref{thm:temp-criterion25}.
Wtedy $\delta_n\to0$ silnie w $\mathcal D'(\R)$:
dla ograniczonego $B\subseteq\mathcal D(\R)$ nośniki
wszystkich jego elementów mieszczą się w jednym zwartym
przedziale (Twierdzenie~\ref{thm:D-bounded18}), więc
$\sup_{\varphi\in B}|\varphi(n)|=0$ dla dużych $n$.
Również $\delta_n\to0$ silnie w $\mathcal S'(\R)$, bo dla
ograniczonego zbioru $B\subseteq\mathcal S$ mamy
$C_B=\sup_{\varphi\in B}p_{2,0}(\varphi)<\infty$ oraz
$\sup_{\varphi\in B}|\varphi(n)|\le C_B(1+n)^{-2}$.
Jednak $\delta_n\not\to0$ nawet słabo-* w
$\mathcal E'(\R)$, ponieważ $\langle\delta_n,1\rangle=1$.

Z kolei $T_n=e^n\delta_n$ zbiega silnie do zera w
$\mathcal D'(\R)$ (na każdym ograniczonym zbiorze
seminorma silna jest od pewnego indeksu zerowa).
Nie zbiega do zera nawet słabo-* w $\mathcal S'(\R)$.
Niech $r(x)=\sqrt{1+x^2}$ i $\varphi(x)=e^{-r(x)/2}$.
Wykażemy, że $\varphi\in\mathcal S(\R)$.
Dla każdego $k\in\N_0$ istnieje wielomian $P_k(x,z)$
taki, że
\[
 \varphi^{(k)}(x)=e^{-r(x)/2}r(x)^{-2k}P_k(x,r(x)).
\]
Dla $k=0$ przyjmujemy $P_0=1$. Ponieważ $r'=x/r$,
różniczkowanie wzoru daje rekurencję
\[
 P_{k+1}=z^2\partial_x P_k+xz\partial_z P_k
          -2kxP_k-\tfrac12 xzP_k.
\]
Zatem wszystkie $P_k$ są wielomianami. Ponieważ
$1\le r(x)\le1+|x|$ i $r(x)\ge|x|$, każda pochodna
jest ograniczona przez stałą razy wielomianowy wzrost
$(1+|x|)^{M_k}$ pomnożony przez $e^{-|x|/2}$.
Wszystkie seminormy Schwartza są więc skończone.
Jednak dla $n\ge1$ mamy $r(n)\le n+1$, toteż
$\langle T_n,\varphi\rangle=e^{n-r(n)/2}
\ge e^{(n-1)/2}\to+\infty$.

Dla dowolnego niepustego zbioru otwartego $\Omega$
wybierzmy wyczerpanie $(K_j)$ i punkty
$a_j\in\Omega\setminus K_j$.
Wówczas $\delta_{a_j}\to0$ silnie w $\mathcal D'(\Omega)$,
gdyż każdy ograniczony zbiór testów ma wspólny zwarty
nośnik zawarty w pewnym $K_J$.
Natomiast $\langle\delta_{a_j},1\rangle=1$ na
$\mathcal E(\Omega)$ dla wszystkich $j$.
Zatem wszystkie odwzorowania z
Wniosku~\ref{cor:dual-lancuch25} nie są izomorfizmami
topologicznymi na swoje obrazy z topologią podprzestrzeni,
ani dla topologii słabych-*, ani dla silnych.
\end{przyklad}

\subsection{Zgodność operacji i granice ich działania}

\begin{stwierdzenie}[Różniczkowanie i mnożenie]
\label{prop:ops25}
Niech $\Omega\subseteq\R^d$ będzie niepustym zbiorem otwartym.
Dla każdego wielowskaźnika $\alpha\in\N_0^d$
operator $\partial^\alpha$ działa liniowo i ciągle na
$\mathcal E(\Omega)$ i $\mathcal D(\Omega)$, a przez
transpozycję także na $\mathcal E'(\Omega)$ i
$\mathcal D'(\Omega)$ w topologiach słabych-* i silnych.
Dla $a\in C^\infty(\Omega)$ mnożenie $f\mapsto af$
i odpowiadające mu mnożenie dystrybucji działają ciągle
na tych przestrzeniach.
Różniczkowanie działa liniowo i ciągle również na
$\mathcal S(\R^d)$ i na $\mathcal S'(\R^d)$ w obu
topologiach dualnych. Na tych przestrzeniach
te same własności mnożenia zachodzą dla
$a\in C^\infty(\R^d)$, jeżeli dla każdego
wielowskaźnika $\alpha$ istnieją $C_\alpha>0$ i
$N_\alpha\in\N_0$ takie, że
$|\partial^\alpha a(x)|\le C_\alpha(1+|x|)^{N_\alpha}$
dla wszystkich $x\in\R^d$.
\end{stwierdzenie}
\begin{proof}
Dla $\mathcal E$ i $\mathcal D$ oszacowania seminorm
dla różniczkowania i mnożenia wynikają z reguły Leibniza,
jak w stwierdzeniu~\ref{prop:operations-cinf17}
i twierdzeniach~\ref{thm:dist-deriv-cont20}, \ref{thm:multiply20}.
Dla $\mathcal S$ i $\mathcal S'$ stosujemy stwierdzenie
\ref{prop:schwartz-operators17} oraz twierdzenie
\ref{thm:tempered-operators22}, w tym oszacowanie
mnożenia przez funkcję o pochodnych wzrostu wielomianowego.
Dla ciągłego operatora $A:X\to X$ transpozycja spełnia
$P_B(A'T)=P_{A(B)}(T)$, gdzie
$P_B(T)=\sup_{f\in B}|T(f)|$ i $A(B)$ jest ograniczony.
Daje to ciągłość silną, a dla singletonów także słabą-*.
Stosujemy ten argument w szczególności do
$X=\mathcal E$ i operatorów $(-1)^{|\alpha|}\partial^\alpha$
oraz $f\mapsto af$. Operatory na funkcjach komutują
z kanonicznymi odwzorowaniami przestrzeni funkcji,
więc przez transpozycję odpowiadające operacje
na przestrzeniach sprzężonych komutują z ograniczeniami
z wniosku~\ref{cor:dual-lancuch25} (w przypadku $\mathcal S'$
przy podanym warunku wzrostu mnożnika).
\end{proof}

\begin{przyklad}[Nie każdy mnożnik gładki zachowuje przestrzeń Schwartza]
\label{ex:not-multiplier25}
Funkcja $a(x)=e^{|x|^2}$ jest gładka, ale mnożenie przez nią
nie zachowuje przestrzeni Schwartza: dla
$f(x)=e^{-|x|^2}\in\mathcal S(\R^d)$ iloczyn $af=1$
nie należy do $\mathcal S(\R^d)$. Przynależność funkcji Gaussa do $\mathcal S$ wynika
z tego, że każda jej pochodna jest wielomianem
pomnożonym przez tę funkcję (indukcja z reguły
iloczynu), a spadek Gaussa dominuje każdą potęgę.
Funkcja $1$ nie spełnia warunku $p_{1,0}(1)<\infty$.
W przestrzeniach $\mathcal D$ i $\mathcal E$
mnożenie przez $a$ jest natomiast ciągłe.
Co więcej, zwykła dystrybucja $aT_1=T_a$ nie jest
temperowana, choć $T_1$ jest temperowana
(Przykłady~\ref{ex:tempered-regular22}
i~\ref{ex:non-tempered22}).
\end{przyklad}

\begin{uwaga}[Transformacja Fouriera i splot]
\label{rem:fourier-conv25}
Nad $\C$ transformacja Fouriera jest automorfizmem
liniowo-topologicznym $\mathcal S(\R^d)$ i
$\mathcal S'(\R^d)$ w obu topologiach dualnych
(Twierdzenia~\ref{thm:fourier-schwartz17}
i~\ref{thm:fourier-auto-tempered22}).
Nie jest zdefiniowana przez tę samą transpozycję
na całej $\mathcal D'$: niezerowa funkcja testowa
nie może mieć transformaty Fouriera o zwartym nośniku.
Istotnie, dla $f\in\mathcal D$ jej transformata
jest rzeczywiście analityczna: rozwinięcie wykładnika
w szereg potęgowy w otoczeniu dowolnego punktu
można całkować wyraz po wyrazie, ponieważ $f$ ma
zwarty nośnik, a szereg jest jednostajnie bezwzględnie
zbieżny na nim dla parametrów z każdego ograniczonego
otoczenia. Gdyby transformata miała zwarty nośnik,
znikałaby na niepustym otwartym podzbiorze $\R^d$,
a z jednoznaczności funkcji rzeczywiście analitycznej
znikałaby na całym spójnym $\R^d$.
Różnowartościowość transformaty na $\mathcal S$
dawałaby $f=0$. Zatem wyrażenie
$\langle T,\mathcal F f\rangle$ nie musi mieć sensu
dla $T\in\mathcal D'$ i $f\in\mathcal D$.

Dla ustalonego $T\in\mathcal E'(\R^d)$ splot
$S\mapsto T*S$ jest ciągłym operatorem liniowym
na $\mathcal D'(\R^d)$ w obu topologiach
(Twierdzenie~\ref{thm:convolution-properties21}).
Aby sprawdzić, że zachowuje również $\mathcal S'$,
niech $L$ będzie zwartym zbiorem i $C>0$, $m\in\N_0$
stałymi oszacowania przedłużenia $T$ na $\mathcal E$
z Twierdzenia~\ref{thm:compact-support-E25}.
Dla $f\in\mathcal S$ określamy
$(A_Tf)(y)=\langle T_x,f(x+y)\rangle$,
korzystając z przedłużenia na $\mathcal E$.
Z ciągłości $T$ na $\mathcal E$ i różniczkowania
przesunięć w jej topologii otrzymujemy
$\partial_y^\beta A_T f(y)=
\langle T_x,\partial^\beta f(x+y)\rangle$:
ilorazy różnicowe wraz z pochodnymi względem $x$
zbiegają jednostajnie na każdym zwartym zbiorze
przez wzór całkowy na różnicę funkcji gładkiej.
Ponieważ $1+|y|\le(1+|x|)(1+|x+y|)$, dla
$R=\max_{x\in L}|x|$ otrzymujemy
\[
 p_{N,k}(A_Tf)\le C(1+R)^N p_{N,k+m}(f).
\]
Zatem $A_T:\mathcal S\to\mathcal S$ jest liniowe
i ciągłe. Jego transpozycja
$\langle T*S,f\rangle=\langle S,A_Tf\rangle$
jest ciągła na $\mathcal S'$ w obu topologiach
i zgadza się na $\mathcal D$ z definicją splotu
z rozdziału~\ref{sec:splot21}.
Nie wynika stąd istnienie splotu dwóch dowolnych
dystrybucji (Uwaga~\ref{rem:convolution-limits21}).
\end{uwaga}

\begin{twierdzenie}[Transformata splotu z dystrybucją o zwartym nośniku]
\label{thm:dt-tempered-compact-convolution}
Niech $U\in\mathcal E'(\R^d)$ oraz $T\in\mathcal S'(\R^d)$.
Definiujemy funkcję
\[
 a(\xi)=c_d\langle U_x,e^{-ix\cdot\xi}\rangle,
 \qquad c_d=(2\pi)^{-d/2},
\]
korzystając z przedłużenia działania $U$ na $\mathcal E$.
Wtedy $a$ jest gładkim mnożnikiem przestrzeni Schwartza,
$\mathcal F U=T_a$, a
\[
 \mathcal F(U*T)=c_d^{-1}a\,\mathcal F T
 \quad\text{w }\mathcal S'.
\]
Splot po lewej stronie jest zgodny ze splotem
dystrybucji z Rozdziału~\ref{sec:splot21}.
\end{twierdzenie}
\begin{proof}
W dalszym rozdziale, w Twierdzeniu~\ref{thm:compact-fourier22},
udowodnimy niezależnie od wzoru na splot, że
$\mathcal FU=T_a$ i że istnieje $m\in\N_0$ takie,
że dla każdego $\alpha\in\N_0^d$ zachodzi
$|\partial^\alpha a(\xi)|\le C_\alpha(1+|\xi|)^m$.
Stosujemy ten wynik tutaj, aby nie powtarzać dowodu
zamiany parowania z całką. Z Twierdzenia
\ref{thm:dt-schwartz-multiplier-characterization}
wynika, że $a$ jest mnożnikiem $\mathcal S$.

Dla ustalonego $y\in\R^d$ i $\varphi\in\mathcal S$
funkcja $\psi_y(\xi)=e^{-iy\cdot\xi}\varphi(\xi)$
należy do $\mathcal S$ na mocy
Twierdzenia~\ref{thm:tempered-operators22}
(jego dowodu dotyczącego mnożników).
Bezpośrednio z definicji całkowej mamy
$(\mathcal F\psi_y)(x)=(\mathcal F\varphi)(x+y)$.
Korzystając z $\mathcal FU=T_a$, dostajemy
\[
 \begin{aligned}
 A_U(\mathcal F\varphi)(y)
 &=\langle U_x,(\mathcal F\varphi)(x+y)\rangle\\
 &=\langle U,\mathcal F\psi_y\rangle
 =\langle\mathcal FU,\psi_y\rangle\\
 &=\int_{\R^d}a(\xi)e^{-iy\cdot\xi}\varphi(\xi)
 \,\mathrm d\lambda_d(\xi)
 =c_d^{-1}\mathcal F(a\varphi)(y),
 \end{aligned}
\]
gdzie $A_Ug(y)=\langle U_x,g(x+y)\rangle$.
Całka jest bezwzględnie zbieżna, ponieważ
$a\varphi\in\mathcal S\subseteq L^1$.
Operator $A_U$ zachowuje $\mathcal S$ i jest ciągły,
jak wykazano w poprzedniej uwadze
\ref{rem:fourier-conv25}. Stąd $U*T$ jest dystrybucją temperowaną.
Korzystamy z definicji splotu przez transpozycję $A_U$:
\[
\begin{aligned}
 \langle\mathcal F(U*T),\varphi\rangle
 &=\langle T,A_U(\mathcal F\varphi)\rangle\\
 &=c_d^{-1}\langle T,\mathcal F(a\varphi)\rangle\\
 &=c_d^{-1}\langle\mathcal F T,a\varphi\rangle .
\end{aligned}
\]
Jest to żądana równość, a zgodność splotów na
$\mathcal D$ wynika z poprzedniej uwagi.
\end{proof}

\begin{uwaga}[Zakres rozdziału]
\label{rem:scope22}
Transformacja Fouriera jest tu automorfizmem przestrzeni
dystrybucji temperowanych i podstawowym narzędziem teorii
operatorów. Jej zastosowania do rozwiązań fundamentalnych rozwiniemy
w dalszych rozdziałach tego skryptu; zastosowania do
konkretnych zagadnień RRC omawia także skrypt o równaniach
różniczkowych cząstkowych.
Dalsze rozdziały rozwiną teorię przestrzeni Montela,
przestrzeni nuklearnych i produktów tensorowych topologicznych.
\end{uwaga}

% Materiał z AF 25; fragment do rozdziału o dystrybucjach temperowanych.
